% Ad Familiares 4.3 --- headnote
% ad-familiares-04-03, August 46 BC (shortly before vi K. intercal. priores), Rome

\textit{Cicero to Servius Sulpicius Rufus, in Achaia, written at
Rome ahead of the sixth day before the earlier intercalary
Kalends of 46 BC --- the Perseus dateline gives \textit{Romae
ante vi K. intercal. priores a. 708 (46)}. Sulpicius, the
distinguished jurist and Cicero's old consular colleague of 51
BC, is now Caesar's governor of Achaia and has been sending
letters home that read --- so Cicero's friends report --- as
the writing of a man too taken up with public misery and his
own absence to take any pleasure in his own goods. Cicero, who
admires Sulpicius's learning and probity above almost any
living contemporary's, writes a consolation modelled on the
Greek philosophical tradition he has been steeping himself in
through the silent years of 46.}

\textit{The letter has three movements that the sections track.
The first builds Sulpicius's claim to clear conscience: he, more
than anyone, foresaw and warned against the civil war from the
consul's chair in 51, and the political ruin that has come is
not his work but the work of those who refused to listen. The
second concedes how thin such consolation feels --- \textit{quasi
parietinis rei publicae}, ``as it were among the ruined walls of
the state'' --- and offers in its place the standing Cicero now
gives Sulpicius in the eyes of Caesar and of every citizen: a
lamp still burning among extinguished lights. The third turns to
philosophy as the remaining shelter, and announces, in passing,
the programme that will produce the philosophical works of the
next two years: ``once I saw that there was no place left for
that art to which I had given my study, neither in the
Senate-house nor in the Forum, I have turned all my care and
labor to philosophy.'' The closing paragraph on the young
Servius --- studying philosophy with Cicero at Rome --- is the
personal warrant for everything that has gone before.}
